
\documentclass{article}
\usepackage{geometry}
\usepackage{fvextra}
\geometry{margin=1in}\usepackage{amsmath}\begin{document}\section*{Fairness Decomposition Report: only LLM }\subsection*{Overview of the Fairness Analysis}This analysis decomposes the disparity in the binary outcome T\_grade between the two sex groups $X=x0$ (F) and $X=x1$ (M). The outcome T\_grade represents whether a student achieved the target grade ($1$ = yes, $0$ = no). We decompose the total group difference into total causal effect, indirect effect through the mediator \texttt{studytime}, direct effect, and spurious effect. There are no confounders ($Z=\text{None}$) and the single ordered mediator is \texttt{studytime}.\subsection*{Decomposition of Effects}\begin{itemize}\item Total variation (observed difference):  \begin{itemize}  \item $P(\text{T\_grade}=1 \mid \text{Sex}= \text{F}) = 0.6394$; $P(\text{T\_grade}=1 \mid \text{Sex}= \text{M}) = 0.7059$.  \item Observed difference (M $-$ F) $= 0.0665$. Interpretation: Males have a $+0.0665$ higher probability of T\_grade $=1$ than Females.  \end{itemize}\item Total causal effect:  \begin{itemize}  \item TE $= E[Y\mid do(X=\text{M})] - E[Y\mid do(X=\text{F})] \approx 0.0665$.  \item Interpretation: Switching Sex from F to M is associated with an increase of about $0.0665$ in the probability of T\_grade $=1$ under the causal assumptions.  \end{itemize}\item Indirect effect (via \texttt{studytime}):  \begin{itemize}  \item Following the baseline convention where mediators use the distribution under $X=x1$ (M) while other components use baseline $X=x0$ (F), the indirect (mediated) effect is    \[      \text{NIE} \;=\; E[Y_{X=\text{F},\;W\sim W_{X=\text{M}}}] - E[Y_{X=\text{F},\;W\sim W_{X=\text{F}}}] \;=\; -0.0012.    \]  \item Interpretation: Replacing the female distribution of \texttt{studytime} by the male distribution (while holding Sex at F) would slightly decrease the probability of success by about $0.0012$ (negligible).  \end{itemize}\item Mediator-specific decomposition:  \begin{itemize}  \item \texttt{studytime} contribution = $-0.0012$ (entire indirect effect).  \item Interpretation: \texttt{studytime} accounts for a negligible and slightly negative share of the disparity.  \end{itemize}\item Direct effect:  \begin{itemize}  \item Direct effect $=$ TE $-$ NIE $= 0.0677$.  \item Interpretation: Most of the positive disparity ($\approx 0.0677$) is attributable to the direct effect of Sex on T\_grade not mediated by \texttt{studytime}.  \end{itemize}\item Spurious effect (via confounders $Z$):  \begin{itemize}  \item No confounders were provided; spurious effect $= 0.0000$.  \item Interpretation: No spurious/confounding contribution is decomposed in this dataset.  \end{itemize}\item X-specific extrema:  \begin{itemize}  \item Direct effect extrema:    \begin{itemize}    \item Largest direct effect: $X=\text{F}$ with direct $= 0.1267$.    \item Smallest direct effect: $X=\text{M}$ with direct $= 0.0677$.    \end{itemize}  \item Indirect effect extrema:    \begin{itemize}    \item Largest indirect effect: $X=\text{F}$ with indirect $= -0.0012$.    \item Smallest indirect effect: $X=\text{M}$ with indirect $= -0.0603$.    \end{itemize}  \item Total effect extrema:    \begin{itemize}    \item Largest total effect: $X=\text{M}$ (effect M vs F) $= +0.0665$.    \item Smallest total effect: $X=\text{F}$ (effect F vs M) $= -0.0665$.    \end{itemize}  \item Interpretation: The direct effect is largest when evaluated with the female baseline; indirect contributions are negligible in magnitude relative to the direct effect.  \end{itemize}\item Z-specific:  \begin{itemize}  \item Z-specific results not provided.  \end{itemize}\end{itemize}\subsection*{Recap}Male students (M) show a higher probability of T\_grade $=1$ than Female students (F) (0.7059 vs 0.6394; difference $= +0.0665$). The decomposition indicates this gap is primarily driven by the direct effect of Sex on T\_grade (direct $= 0.0677$), while the mediator \texttt{studytime} contributes negligibly and slightly negatively (indirect $= -0.0012$), and there is no spurious/confounding contribution in the provided data. In substantive terms, the observed advantage for M over F in T\_grade in this dataset is mainly attributable to direct pathways rather than studytime or measured confounding.

\begin{Verbatim}[breaklines=true]
 ResponseUsage(input_tokens=7905, input_tokens_details=InputTokensDetails(cached_tokens=1792), output_tokens=7245, output_tokens_details=OutputTokensDetails(reasoning_tokens=4544), total_tokens=15150)
\end{Verbatim}
\end{document}